\subsection{Finitely presented modules}

Let \(\mathbf{A}\) be a ring. An exact sequence

\[
\mathrm{L} _ {1} \rightarrow \mathrm{L} _ {0} \rightarrow \mathrm{E} \rightarrow 0\tag{6}
\]

of left (resp. right) A-modules, where \(L_{0}\) and \(L_{1}\) are free, is called a presentation (or presentation of length 1) of a left (resp. right) A-module E.

Every A-module E admits a presentation. We know in fact (Algebra, Chapter II, § 1, no. 11, Proposition 20) that there exists a surjective homomorphism \(u \colon \mathbf{L}_0 \to \mathbf{E}\) , where \(\mathbf{L}_0\) is free; if \(\mathbf{R}\) is the kernel of \(u\) , there exists similarly a surjective homomorphism \(v \colon \mathbf{L}_1 \to \mathbf{R}\) , where \(\mathbf{L}_1\) is free. If \(v\) is considered as a homomorphism of \(\mathbf{L}_1\) to \(\mathbf{L}_0\) , the sequence \(\mathbf{L}_1 \xrightarrow{v} \mathbf{L}_0 \xrightarrow{u} 0\) is exact by definition, whence our assertion.

If \(\rho: \mathbf{A} \to \mathbf{B}\) is a ring homomorphism, every presentation (6) of \(\mathbf{E}\) induces a presentation of \(\mathbf{E}_{(\mathbf{B})} = \mathbf{E} \otimes_{\mathbf{A}} \mathbf{B}\) :

\[
\mathrm{L} _ {1} \otimes_ {\mathrm{A}} \mathrm{B} \rightarrow \mathrm{L} _ {0} \otimes_ {\mathrm{A}} \mathrm{B} \rightarrow \mathrm{E} \otimes_ {\mathrm{A}} \mathrm{B} \rightarrow 0\tag{7}
\]

by Lemma 1 of no. 1 and the fact that \(\mathbf{L} \otimes_{\mathbf{A}} \mathbf{B}\) is a free B-module if \(\mathbf{L}\) is free.

A presentation (6) of a module E is called finite if the free modules \(L_{0}\) and \(L_{1}\) have finite bases. Clearly, if the presentation (6) is finite, so is the presentation (7). E is called a finitely presented A-module if it admits a finite presentation.

Assertion (i) follows trivially from the definitions. If A is left Noetherian and there exists a surjective homomorphism \(u \colon \mathbf{L}_0 \to \mathbf{E}\) , where \(\mathbf{L}_0\) is a free left A-module with a finite basis, the kernel R of \(u\) is finitely generated (Algebra, Chapter VIII, § 2, no. 1, Proposition 1 and no. 3, Proposition 7), hence there is a surjective homomorphism \(v \colon \mathbf{L}_1 \to \mathbf{R}\) , where \(\mathbf{L}_1\) is free and has a finite basis, and the exact sequence \(\mathbf{L}_1 \xrightarrow{v} \mathbf{L}_0 \xrightarrow{u} \mathbf{E} \to 0\) is a finite presentation of E; whence (ii).

Finally, suppose that E is a finitely generated projective module; then it is a direct factor of a finitely generated free module \(L_{0}\) (Algebra, Chapter II, §2, no. 2, Corollary to Proposition 4); the kernel R of the surjective homomorphism \(L_{0} \rightarrow E\) is then isomorphic to a quotient of \(L_{0}\) and hence finitely generated and the proof is completed as for (ii).

LEMMA 9. Let A be a ring and E a finitely presented A-module. For every exact sequence

\[
0 \longrightarrow \mathrm{F} \stackrel {j} {\longrightarrow} \mathrm{G} \stackrel {p} {\longrightarrow} \mathrm{E} \longrightarrow 0
\]

where G is finitely generated, the module F is finitely generated.

Let \(L_{1} \stackrel{r}{\rightarrow} L_{0} \stackrel{s}{\rightarrow} E \rightarrow 0\) be a finite presentation; if \((e_{i})\) is a basis of \(L_{0}\) , there exists for each i an element \(g_{i} \in G\) such that \(p(g_{i}) = s(e_{i})\) ; the homomorphism \(u: L_{0} \rightarrow G\) such that \(u(e_{i}) = g_{i}\) for all i then satisfies \(s = p \circ u\) . As \(s \circ r = 0\) , we have \(u(r(L_{1})) \subset \operatorname{Ker} p\) , and as \(\operatorname{Ker} p\) is isomorphic to F, it can be seen that there is a homomorphism \(v: L_{1} \rightarrow F\) such that the diagram

\[
\begin{array}{c} \mathrm{L} _ {1} \xrightarrow {r} \mathrm{L} _ {0} \xrightarrow {s} \mathrm{E} \longrightarrow 0 \\ v \Biggl \downarrow \qquad \qquad u \Biggl \downarrow \qquad \qquad 1 _ {\mathrm{E}} \Biggl \downarrow \\ \mathrm{F} \xrightarrow {} _ {j} \mathrm{G} \xrightarrow {} _ {p} \mathrm{E} \longrightarrow 0 \end{array}
\]

is commutative. As \(j\) is injective and \(s\) surjective, we can apply the snake diagram (§ 1, no. 4, Proposition 4), in other words there is an exact sequence

\[
0 = \operatorname{Ker} 1 _ {\mathrm{E}} \xrightarrow {d} \text { Coker } v \longrightarrow \text { Coker } u \longrightarrow \text { Coker } 1 _ {\mathrm{E}} = 0.
\]

This shows that Coker v is isomorphic to \(G/u(L_{0})\) , which is finitely generated by hypothesis. Moreover we have the exact sequence

\[
0 \rightarrow v (\mathrm{L} _ {1}) \rightarrow \mathrm{F} \rightarrow \text { Coker } v \rightarrow 0
\]

and as \(v(\mathbf{L}_1)\) and Coker \(v\) are finitely generated, so is F (Algebra, Chapter II, § 1, no. 7, Corollary 5 to Proposition 9).

